\section{Dice instantiation and relation to SDC}
\label{sec:theory}

For binary masks, define
\begin{equation}
  \Dice(A,B)=\frac{2|A\cap B|}{|A|+|B|},
  \label{eq:dice-definition}
\end{equation}
with \(\Dice(\varnothing,\varnothing)=1\) and Dice equal to zero when
exactly one mask is empty. The loss \(L_{\Dice}=1-\Dice\) lies in
\([0,1]\). Equation~\eqref{eq:quad} gives
\begin{align}
  C_{\Dice,M}(x)
  &:=-\sum_{m=1}^M w_m
      \left[1-\Dice(Y_{t_m},\Yhat)\right], \\
  1+C_{\Dice,M}(x)
  &=\sum_{m=1}^M w_m\Dice(Y_{t_m},\Yhat).
  \label{eq:levelset-dice}
\end{align}
Thus \(1+C_{\Dice,M}\) estimates expected Dice under \(Q_p\) and is
rank-equivalent to the loss-oriented confidence.

When the denominator is positive, SDC computes the soft Dice between \(p\)
and the deployed mask:
\begin{equation}
  \SDC(p,\Yhat)
  =
  \frac{2\sum_i p_i\hat y_i}
       {\sum_i p_i+\sum_i\hat y_i},
  \qquad
  \hat y_i=\mathbf 1\{i\in\Yhat\}.
  \label{eq:sdc}
\end{equation}
The published SDC convention assigns zero when both arguments vanish
\citep{sdc}. This differs from our deployment-loss convention
\(\Dice(\varnothing,\varnothing)=1\), a distinction that matters when the
deployed mask is empty. It does not affect the following identity, which
assumes a positive denominator.

\begin{proposition}[SDC is a ratio of expected Dice components]
\label{prop:sdc}
Let \(Q\) be any distribution over binary masks with marginals
\(\Ex_Q[Y_i]=p_i\), and let
\(N(Y)=2\sum_iY_i\hat y_i\) and
\(D(Y)=\sum_iY_i+\sum_i\hat y_i\). If \(\Ex_Q[D(Y)]>0\), then
\begin{equation}
  \frac{\Ex_Q[N(Y)]}{\Ex_Q[D(Y)]}
  =\SDC(p,\Yhat).
  \label{eq:sdc-ratio-expectations}
\end{equation}
The right-hand side is therefore invariant to the coupling among pixels once
the marginals are fixed.
\end{proposition}
Both \(N\) and \(D\) are affine in \(Y\), so the identity follows by
linearity of expectation. If \(D(Y)>0\) almost surely---in particular, whenever
\(\Yhat\neq\varnothing\)---expected Dice is an expectation of a ratio:
\begin{equation}
  \Ex_Q[\Dice(Y,\Yhat)]
  =\Ex_Q\!\left[\frac{N(Y)}{D(Y)}\right]
  \neq
  \frac{\Ex_Q[N(Y)]}{\Ex_Q[D(Y)]}
  \quad\text{in general}.
  \label{eq:ratio-noncommute}
\end{equation}
When \(D(Y)=0\) has positive probability, the left-hand side instead includes
the declared empty--empty value on that event; it still is not determined by
the ratio of the two expectations.

A two-pixel example makes the distinction concrete. Let
\(p=(1/2,1/2)\) and \(\Yhat=\{1\}\). Under independent Bernoulli pixels,
\(\Ex[\Dice(Y,\Yhat)]=5/12\). Under the level-set coupling, \(Y\) is
either \(\{1,2\}\) or \(\varnothing\), each with probability \(1/2\), so
\(\Ex_{Q_p}[\Dice(Y,\Yhat)]=1/3\). Both couplings have the same marginals
and \(\SDC(p,\Yhat)=1/2\). The three quantities are distinct.

\paragraph{Assumptions are not ordered.}
Computing SDC requires no explicit joint posterior. Interpreting it as an
approximation to true ideal expected Dice requires the input map to represent
the relevant true marginals and, in the published analysis, conditional pixel
independence. Under the product law induced by those marginals, the SDC theory
directly bounds the ratio between SDC and its marginal-based ideal Dice
confidence \citep{sdc}. Because that analysis sets
\(\Dice(\varnothing,\varnothing)=0\), its guarantee transfers directly to our
Dice convention for nonempty deployed masks, but not for an empty deployed
mask without an additional argument.
The level-set estimator instead evaluates expected Dice exactly under its
explicit comonotone working posterior in the dense-quadrature limit. Its link
to the true ideal confidence requires the loss-specific discrepancy in
Equation~\eqref{eq:loss-discrepancy} to be small.

The structural conditions are not ordered. Conditional independence does not
force exact agreement with \(Q_p\): the two-pixel example above has independent
true labels but Dice discrepancy \(|5/12-1/3|=1/12\) from the level-set law. Conversely,
labels generated by a shared threshold make \(Q_p\) exact while violating
conditional independence except in degenerate cases. Approximate conditions
must therefore be compared at explicit tolerances rather than labeled weaker
in the abstract. The framework generalizes the
\emph{loss-indexed principle} beyond Dice, but does not contain SDC as a literal
special case or claim uniformly weaker probabilistic assumptions.
